\documentclass[11pt,a4paper]{article}
\usepackage[margin=2.5cm]{geometry}
\usepackage{amsmath,amssymb}
\usepackage{graphicx}
\usepackage{booktabs}
\usepackage{longtable}
\usepackage{array}
\usepackage{xcolor}
\usepackage{hyperref}
\usepackage{listings}
\usepackage{fancyhdr}
\usepackage{titlesec}
\usepackage{enumitem}
\usepackage{mdframed}

\definecolor{sambpblue}{RGB}{0,70,127}
\definecolor{codegreen}{RGB}{0,128,0}
\definecolor{codegray}{RGB}{128,128,128}
\definecolor{codeblue}{RGB}{0,0,200}
\definecolor{codebg}{RGB}{248,248,248}
\definecolor{lightblue}{RGB}{230,240,255}
\definecolor{lightyellow}{RGB}{255,253,230}

\hypersetup{
  colorlinks=true,
  linkcolor=sambpblue,
  urlcolor=sambpblue,
  citecolor=sambpblue
}

\lstdefinestyle{python}{
  language=Python,
  backgroundcolor=\color{codebg},
  basicstyle=\ttfamily\footnotesize,
  keywordstyle=\color{codeblue}\bfseries,
  commentstyle=\color{codegreen}\itshape,
  stringstyle=\color{orange},
  numberstyle=\tiny\color{codegray},
  numbers=left,
  numbersep=8pt,
  breaklines=true,
  frame=single,
  framerule=0.4pt,
  rulecolor=\color{gray!40},
  xleftmargin=1.5em,
  framexleftmargin=1.5em,
  captionpos=b,
  showstringspaces=false,
  tabsize=4,
}

\lstdefinestyle{bash}{
  language=bash,
  backgroundcolor=\color{black!90},
  basicstyle=\ttfamily\footnotesize\color{green!70!black},
  keywordstyle=\color{yellow},
  commentstyle=\color{green!50},
  stringstyle=\color{orange!80},
  frame=none,
  xleftmargin=1em,
  breaklines=true,
  showstringspaces=false,
}

\pagestyle{fancy}
\fancyhf{}
\fancyhead[L]{\small\textcolor{sambpblue}{\textbf{SAMBP -- Simulation Tool Integration Report}}}
\fancyhead[R]{\small\textcolor{sambpblue}{IIT Madras PhD Thesis}}
\fancyfoot[C]{\small\thepage}
\fancyfoot[L]{\small\textcolor{gray}{Anoop V Eluvathingal}}
\fancyfoot[R]{\small\textcolor{gray}{April 2026}}
\renewcommand{\headrulewidth}{0.5pt}
\renewcommand{\footrulewidth}{0.3pt}

\titleformat{\section}{\large\bfseries\color{sambpblue}}{{\thesection}}{0.8em}{}[\vspace{-0.3em}\rule{\linewidth}{0.5pt}]
\titleformat{\subsection}{\normalsize\bfseries\color{sambpblue!80}}{\thesubsection}{0.7em}{}
\titleformat{\subsubsection}{\normalsize\bfseries\color{black}}{\thesubsubsection}{0.6em}{}

\newmdenv[backgroundcolor=lightblue,linecolor=sambpblue,linewidth=1pt,
          leftmargin=0,rightmargin=0,innertopmargin=6pt,innerbottommargin=6pt,
          innerleftmargin=10pt,innerrightmargin=10pt]{infobox}

\newmdenv[backgroundcolor=lightyellow,linecolor=orange!60,linewidth=1pt,
          leftmargin=0,rightmargin=0,innertopmargin=6pt,innerbottommargin=6pt,
          innerleftmargin=10pt,innerrightmargin=10pt]{notebox}

\begin{document}

% ============================================================
%  TITLE PAGE
% ============================================================
\begin{titlepage}
\centering
\vspace*{1cm}
{\color{sambpblue}\rule{\linewidth}{2pt}}\\[0.4cm]
{\huge\bfseries\color{sambpblue} Power System Simulation Tool Integration\\[0.3cm]
for SAMBP Research}\\[0.3cm]
{\color{sambpblue}\rule{\linewidth}{1pt}}\\[1cm]

{\large\textit{Integrating ANDES, OpenDSS, and pandapower\\
with the Self-Adaptive Model-Based Protection Framework}}\\[2cm]

\begin{tabular}{rl}
\textbf{Author}    & Anoop V Eluvathingal \\[4pt]
\textbf{Degree}    & Doctor of Philosophy, Electrical Engineering \\[4pt]
\textbf{Institute} & IIT Madras \\[4pt]
\textbf{Date}      & April 2026 \\[4pt]
\textbf{Version}   & 1.0 -- Initial Analysis
\end{tabular}\\[2cm]

\begin{infobox}
\textbf{Document Purpose}\\[4pt]
This report analyses how to integrate open-source power system dynamic
simulation tools (ANDES, OpenDSS, pandapower, Dynawo, GridPACK) with the
SAMBP protection framework running on the VPS. It provides:
(1)~a tool selection rationale, (2)~a module-by-module integration
architecture, (3)~ready-to-use Python adapter code, and (4)~a concrete
installation and execution plan.
\end{infobox}

\vfill
{\color{sambpblue}\rule{\linewidth}{1pt}}
\end{titlepage}

\tableofcontents
\newpage

% ============================================================
\section{VPS Environment and Available Tools}
% ============================================================

\subsection{VPS Hardware Specifications}
The current VPS on which SAMBP runs has the following capabilities:

\begin{center}
\begin{tabular}{ll}
\toprule
\textbf{Parameter} & \textbf{Value} \\
\midrule
Architecture & x86\_64 \\
RAM          & 15 GB \\
Disk         & 193 GB (25 GB used, 169 GB free) \\
Python       & 3.12.3 \\
OS           & Linux (Ubuntu) \\
\bottomrule
\end{tabular}
\end{center}

This is sufficient to run ANDES transient simulations on networks up to
the IEEE 39-bus scale, pandapower Monte Carlo sweeps with thousands of
cases, and OpenDSS distribution studies.

\subsection{Tools Shown in the Slide}

The NREL slide identifies four active open-source simulation ecosystems:

\begin{center}
\begin{tabular}{llll}
\toprule
\textbf{Tool} & \textbf{Developer} & \textbf{Language} & \textbf{Focus} \\
\midrule
ANDES             & CURENT / Texas A\&M & Python            & Transmission transient dynamics \\
Sienna$\backslash$Dyn & NREL           & Julia             & Transmission transient dynamics \\
Dynawo            & RTE + LF Energy     & C++/Modelica      & Transmission + distribution \\
GridPACK          & Pacific NW Lab      & C++ (MPI)         & HPC parallel dynamics \\
\bottomrule
\end{tabular}
\end{center}

\begin{notebox}
\textbf{Not on the slide but equally important for SAMBP:}
\textbf{pandapower} (Fraunhofer IEE) and \textbf{OpenDSS / opendssdirect}
(EPRI) -- both pure Python, directly pip-installable, and covering
distribution-level protection studies that ANDES does not target.
\end{notebox}

% ============================================================
\section{Tool Selection Rationale}
% ============================================================

\subsection{Selection Criteria for SAMBP}

SAMBP needs a simulation back-end that can:
\begin{enumerate}[noitemsep]
\item Produce time-domain three-phase current waveforms \(i_a(t), i_b(t), i_c(t)\) at relay terminal buses
\item Model IBR current-limiting behaviour (GFL/GFM converters, DFIG, PV)
\item Allow parametric fault sweeps (location, type, resistance, IBR penetration)
\item Integrate with Python without a separate GUI or licence
\item Run on the VPS without MPI/HPC infrastructure
\end{enumerate}

\subsection{Comparative Analysis}

\begin{longtable}{p{2.2cm}p{2.4cm}p{1.8cm}p{2.8cm}p{2.8cm}}
\toprule
\textbf{Tool} & \textbf{Simulation type} & \textbf{Python API} & \textbf{IBR support} & \textbf{Fit for SAMBP} \\
\midrule
\endhead
\textbf{ANDES}
  & Transmission transient (DAE, RK4)
  & Native Python
  & REGCA1 (GFL), DFIG, PV plant (REECA1/REPCA1)
  & \textbf{Best -- covers SG + IBR transient for 87T, 87L, OC, Relay 78} \\[6pt]
\textbf{OpenDSS / opendssdirect}
  & Distribution quasi-static + dynamic
  & Python bindings
  & Static IBR (PQ model)
  & \textbf{Best for 87B, 87L distribution, HIF (TR-11)} \\[6pt]
\textbf{pandapower}
  & Power flow + short circuit
  & Native Python
  & Static IBR (PQ)
  & \textbf{Good for MC sweeps, relay coordination setting studies} \\[6pt]
\textbf{Sienna$\backslash$Dyn}
  & Transmission transient
  & Julia only
  & Full IBR models
  & Not practical -- requires Julia runtime, no Python API \\[6pt]
\textbf{Dynawo}
  & Transmission + distribution
  & C++ / Python co-sim
  & Modelica IBR models
  & Overkill -- 500 MB install, Modelica learning curve \\[6pt]
\textbf{GridPACK}
  & HPC parallel dynamics
  & C++ / Python
  & Yes
  & Overkill -- requires MPI build, not needed at thesis scale \\
\bottomrule
\end{longtable}

\begin{infobox}
\textbf{Recommended stack for VPS:}\\[4pt]
\texttt{pip install andes pandapower opendssdirect}\\[4pt]
This covers all nine SAMBP protection functions with no additional
infrastructure. Estimated install size: $\approx$100 MB combined.
\end{infobox}

% ============================================================
\section{Integration Architecture}
% ============================================================

\subsection{System Overview}

The SAMBP framework currently generates waveforms through two synthetic
models: \texttt{network\_fault\_model.py} (Th\'evenin equivalent, per
time-step flags) and \texttt{radial\_network.py} (analytical three-bus
network). The integration replaces only that data-generation layer with
physics-correct simulation output while leaving all downstream SAMBP
logic \textit{unchanged}.

\begin{center}
\renewcommand{\arraystretch}{1.4}
\begin{tabular}{|c|}
\hline
\textbf{LAYER 1 -- Simulation}\\
ANDES / OpenDSS / pandapower\\
Define topology $\to$ inject fault $\to$ extract $I_{abc}(t)$, $V_{abc}(t)$\\
\hline
$\downarrow$ \quad CSV / NumPy arrays \quad $\downarrow$ \\
\hline
\textbf{LAYER 2 -- SAMBP Data Adapter} (new)\\
\texttt{04\_code/sambp/io\_utils/andes\_adapter.py}\\
\texttt{04\_code/sambp/io\_utils/opendss\_adapter.py}\\
Resample to 10 kHz $\cdot$ extract per-relay terminals $\cdot$ package FaultCase\\
\hline
$\downarrow$ \quad FaultCase dataclass \quad $\downarrow$ \\
\hline
\textbf{LAYER 3 -- SAMBP Core} (unchanged)\\
\texttt{inverse\_estimation/} $\to$ \texttt{adaptation/} $\to$ Stage-2 gate\\
\hline
\end{tabular}
\end{center}

\subsection{Parameter Mapping: SAMBP to ANDES}

Your existing \texttt{system\_config.py} Park dq parameters map
directly to the ANDES \texttt{GENROU} device:

\begin{center}
\begin{tabular}{lll}
\toprule
\textbf{SAMBP \texttt{system\_config.py}} & \textbf{ANDES GENROU} & \textbf{Description} \\
\midrule
\texttt{Xd\_pu}   & \texttt{xd}    & d-axis synchronous reactance \\
\texttt{Xdp\_pu}  & \texttt{xd1}   & d-axis transient reactance \\
\texttt{Xdpp\_pu} & \texttt{xd2}   & d-axis subtransient reactance \\
\texttt{Xq\_pu}   & \texttt{xq}    & q-axis synchronous reactance \\
\texttt{Xqp\_pu}  & \texttt{xq1}   & q-axis transient reactance \\
\texttt{Tdp\_s}   & \texttt{Td10}  & d-axis open-circuit transient TC \\
\texttt{Tdpp\_s}  & \texttt{Td20}  & d-axis subtransient TC \\
\texttt{H\_s}     & \texttt{H}     & Inertia constant (MW$\cdot$s/MVA) \\
\texttt{D\_pu}    & \texttt{D}     & Damping coefficient \\
\texttt{Ka, Ta\_s}& \texttt{ESST1A}& AVR gain and time constant \\
\bottomrule
\end{tabular}
\end{center}

% ============================================================
\section{Module-by-Module Integration}
% ============================================================

\subsection{sync\_oc -- ANDES Transient Simulation}

The \texttt{sync\_oc} module estimates the synchronous generator
parameters ($X_d''$, $R_a$, $\tau$) from fault current waveforms.
Currently it uses a synthetic waveform from \texttt{network\_fault\_model.py}.
Replacing with ANDES:

\begin{lstlisting}[style=python, caption={04\_code/sambp/io\_utils/andes\_adapter.py}]
import andes
import numpy as np
from scipy.interpolate import interp1d

def run_andes_fault(case_file: str,
                    fault_bus: int,
                    fault_time: float,
                    clear_time: float,
                    fs_target: float = 10000.0) -> dict:
    """
    Run ANDES transient simulation and return waveforms resampled
    to SAMBP sampling rate (default 10 kHz).
    """
    ss = andes.load(case_file)
    ss.PFlow.run()                        # initialise power flow

    # Add fault event via ANDES Fault device
    ss.add("Fault", {
        "bus": fault_bus,
        "tf":  fault_time,
        "tc":  clear_time,
        "xf":  0.01          # fault reactance [pu]
    })

    ss.TDS.config.tf = clear_time + 0.30  # 300 ms post-fault window
    ss.TDS.run()

    # Extract time vector and bus voltages
    t = np.array(ss.dae.ts.t)
    V = np.array(ss.dae.ts.y)[:, ss.Bus.v.a]   # all bus voltages [pu]

    # Resample to SAMBP fs_target
    t_new = np.arange(t[0], t[-1], 1.0 / fs_target)
    V_rs  = interp1d(t, V, axis=0)(t_new)

    return {"t": t_new, "V_pu": V_rs, "ss": ss}


def get_branch_current(ss, from_bus: int, to_bus: int) -> np.ndarray:
    """Extract complex current phasor time-series on branch from->to."""
    # ANDES Line model stores Ixy in dae.ts
    line_idx = _find_line(ss, from_bus, to_bus)
    I_re = np.array(ss.dae.ts.y)[:, ss.Line.a[line_idx]]
    I_im = np.array(ss.dae.ts.y)[:, ss.Line.v[line_idx]]
    return I_re + 1j * I_im


def build_sambp_case(t, I_near, I_far, V_near, V_far,
                     fault_type: str, fault_R: float) -> dict:
    """Package ANDES output as SAMBP FaultCase dictionary."""
    return {
        "t":          t,
        "I_near_pu":  I_near,
        "I_far_pu":   I_far,
        "V_near_pu":  V_near,
        "V_far_pu":   V_far,
        "fault_type": fault_type,
        "fault_R_pu": fault_R,
    }
\end{lstlisting}

Usage in the existing \texttt{sync\_oc} pipeline:
\begin{lstlisting}[style=python, caption={Run sync\_oc with ANDES instead of synthetic data}]
from io_utils.andes_adapter import run_andes_fault, build_sambp_case
from inverse_estimation.lm_estimator import run_two_pass_lm

# Previously: data from network_fault_model.py (synthetic Thevenin)
# Now:        data from ANDES transient simulation
sim = run_andes_fault("ieee9.json", fault_bus=5,
                       fault_time=0.10, clear_time=0.25)

case = build_sambp_case(
    t       = sim["t"],
    I_near  = sim["I_near"],
    I_far   = sim["I_far"],
    V_near  = sim["V_near"],
    V_far   = sim["V_far"],
    fault_type = "SLG",
    fault_R    = 0.01,
)
result = run_two_pass_lm(case)   # unchanged -- no modifications needed
\end{lstlisting}

\subsection{line\_diff (87L) -- ANDES Line Models}

For the 87L protection (TR-12--TR-27, TR-56 DFIG, TR-62 PV), ANDES
provides full transient current at both line ends. The two-pass LM
estimator in \texttt{line\_diff/inverse\_estimation/} consumes exactly
these near-end and far-end complex current phasors.

For IBR-connected lines, replace \texttt{GENROU} with the ANDES
renewable energy model chain:

\begin{lstlisting}[style=python, caption={Adding DFIG/PV in ANDES for TR-56/TR-62}]
# DFIG (TR-56: piecewise-ODE model)
ss.add("DFIG", {
    "bus":    gen_bus,
    "Sn":     100.0,       # MVA rating
    "Vn":     0.69,        # kV
    "fn":     50.0,
    "rs":     0.01,
    "Lls":    0.10,
    "Lm":     3.50,
    "Llr":    0.09,
    "rr":     0.01,
    "H":      5.04,
})

# PV plant (TR-62: 2-parameter zone model)
ss.add("REGCA1", {
    "bus":    pv_bus,
    "Tg":     0.02,
    "Rrpwr":  10.0,
    "Brkpt":  0.90,
    "Zerox":  0.40,
    "Lvpl1":  1.22,
    "Iqmax":  1.10,        # current limit [pu] -- key IBR parameter
})
ss.add("REECA1", {
    "reg":    pv_idx,
    "PQFLAG": 0,           # Q-priority during fault
    "Vdip":   0.90,
    "Vup":    1.10,
    "Tiq":    0.02,
    "dPmax":  3.0,
    "dPmin": -3.0,
    "Pmax":   1.0,
    "Iqh1":   1.10,
    "Iql1":  -1.10,
})
\end{lstlisting}

\subsection{transformer\_diff (87T) -- ANDES Transformer}

The ANDES \texttt{Transformer} device gives HV and LV terminal current
vectors directly. No changes to \texttt{transformer\_diff/models/} are
needed:

\begin{lstlisting}[style=python, caption={Extract transformer terminal currents}]
def get_transformer_current(ss, tr_idx: int, side: str) -> np.ndarray:
    """
    side: 'hv' or 'lv'
    Returns complex current time-series at transformer terminal [pu].
    """
    tr = ss.Transformer
    if side == "hv":
        bus_idx = tr.bus1.v[tr_idx]
    else:
        bus_idx = tr.bus2.v[tr_idx]
    # Build from dae time-series (voltage / impedance)
    V = np.array(ss.dae.ts.y)[:, bus_idx]
    Z = complex(tr.r.v[tr_idx], tr.x.v[tr_idx])
    return V / Z   # Ohm's law in phasor domain
\end{lstlisting}

\subsection{bus\_diff (87B) -- OpenDSS for Distribution Buses}

OpenDSS handles multi-feeder distribution buses better than ANDES.
The \texttt{opendssdirect} Python API feeds directly into
\texttt{bus\_diff/models/bus\_reduced\_zone\_model.py}:

\begin{lstlisting}[style=python, caption={04\_code/sambp/io\_utils/opendss\_adapter.py}]
import opendssdirect as dss
import numpy as np

def run_opendss_fault(dss_file: str,
                      fault_bus: str,
                      fault_R: float = 0.01,
                      n_phases: int = 3,
                      t_step: float = 1e-4,
                      n_steps: int = 2000) -> dict:
    """
    Run OpenDSS dynamic simulation with a fault on fault_bus.
    Returns feeder currents at each element connected to fault_bus.
    """
    dss.run_command(f"Compile {dss_file}")
    dss.run_command(
        f"New Fault.F1 Bus1={fault_bus} Phases={n_phases} r={fault_R}"
    )
    dss.run_command(
        f"Set Mode=Dynamic StepSize={t_step} Number={n_steps}"
    )
    dss.Solution.Solve()

    # Extract currents at each element feeding fault_bus
    I_feeders = {}
    for elem in dss.Circuit.AllElementNames():
        dss.Circuit.SetActiveElement(elem)
        buses = dss.CktElement.BusNames()
        if any(fault_bus in b for b in buses):
            I_raw = dss.CktElement.Currents()   # [Re, Im, Re, Im, ...]
            I_feeders[elem] = (
                np.array(I_raw[0::2]) + 1j * np.array(I_raw[1::2])
            )

    t = np.arange(n_steps) * t_step
    return {"t": t, "I_feeders": I_feeders}


def package_for_bus_diff(sim: dict, feeder_names: list) -> dict:
    """
    Map OpenDSS feeder currents to the bus_diff zone model input format.
    feeder_names: ordered list matching bus_reduced_zone_model expectations.
    """
    return {
        "t":         sim["t"],
        "I_feeders": [sim["I_feeders"].get(n, np.zeros(len(sim["t"])))
                      for n in feeder_names],
    }
\end{lstlisting}

\subsection{Monte Carlo Validation -- pandapower}

Replace the synthetic inner loop in \texttt{sync\_oc/main\_batch\_study.py}
with pandapower parametric sweeps to generate
physics-backed MC validation data:

\begin{lstlisting}[style=python, caption={MC sweep using pandapower (upgrade to main\_batch\_study.py)}]
import pandapower as pp
import pandapower.networks as pn
import numpy as np
import pandas as pd
from sync_oc.inverse_estimation.lm_estimator import run_two_pass_lm

def set_ibr_penetration(net, ibr_frac: float):
    """Replace ibr_frac of total load with IBR generation."""
    total_load = net.load.p_mw.sum()
    ibr_mw = ibr_frac * total_load
    # Add IBR as negative load (constant PQ) at bus 3
    net.load.at[0, "p_mw"] -= ibr_mw

results = []
net_base = pn.case33bw()   # IEEE 33-bus distribution network

for fault_R in np.linspace(0.01, 50.0, 100):       # fault resistance [Ohm]
    for ibr_pen in [0.0, 0.3, 0.5, 0.8, 1.0]:     # IBR penetration ratio
        net = net_base.deepcopy()
        set_ibr_penetration(net, ibr_pen)
        pp.runpp(net, algorithm="nr")

        # Compute three-phase fault current at bus 5
        pp.shortcircuit.calc_sc(net, bus=5, fault="3ph",
                                 r_fault_ohm=fault_R)
        I_fault_kA = net.res_bus_sc.at[5, "ikss_ka"]

        # Feed into SAMBP estimator
        I_wave = _synthetic_wave_from_phasor(I_fault_kA, t_end=0.3)
        result = run_two_pass_lm({"I_near_pu": I_wave, ...})
        results.append({
            "fault_R_ohm":  fault_R,
            "ibr_pen":      ibr_pen,
            "I_fault_kA":   I_fault_kA,
            "kappa_est":    result["kappa"],
            "stage2_pass":  result["stage2_gate"],
        })

df = pd.DataFrame(results)
df.to_csv("05_data/mc_pandapower_results.csv", index=False)
\end{lstlisting}

% ============================================================
\section{IBR-Specific Considerations}
% ============================================================

\subsection{Why IBR Changes Everything for Protection}

The fundamental challenge that SAMBP addresses is that IBRs do not
behave like synchronous generators during faults:

\begin{center}
\begin{tabular}{lcc}
\toprule
\textbf{Characteristic} & \textbf{SG (classic)} & \textbf{IBR (GFL)} \\
\midrule
Fault current magnitude & 5--8 $\times$ rated & 1.05--1.20 $\times$ rated \\
Current waveform        & Decaying exponential  & Controlled, near-sinusoidal \\
Phase angle change      & Large (lagging)       & Small (regulator-set) \\
Recovery time           & Governed by $X_d''$   & Governed by PLL bandwidth \\
Zero-sequence path      & Available             & Absent (converter-isolated) \\
\bottomrule
\end{tabular}
\end{center}

ANDES reproduces this behaviour faithfully through its \texttt{REGCA1}
(generator-converter) model, which enforces the \texttt{Iqmax} current
ceiling automatically. The SAMBP Stage-2 veto gate uses exactly this
reduced current signature to distinguish IBR from SG -- making ANDES
the ideal physics engine to validate Stage-2 performance.

\subsection{ANDES IBR Model Chain for TR-56/TR-62}

\begin{center}
\begin{tabular}{llll}
\toprule
\textbf{ANDES model} & \textbf{Layer} & \textbf{SAMBP TR} & \textbf{Function} \\
\midrule
\texttt{REGCA1} & Converter dynamics & TR-62 (PV) & Current limiter, LVRT \\
\texttt{REECA1} & Electrical control & TR-56, TR-62 & P/Q priority switching \\
\texttt{REPCA1} & Plant control      & TR-62 (PV plant) & AGC, volt-var control \\
\texttt{DFIG}   & DFIG machine       & TR-56            & Rotor + stator dynamics \\
\texttt{ESST1A} & AVR (for SG)       & TR-57 (87T IBR)  & Exciter dynamics \\
\bottomrule
\end{tabular}
\end{center}

% ============================================================
\section{Installation and Quick-Start Guide}
% ============================================================

\subsection{Step 1: Install the Tool Stack}

\begin{lstlisting}[style=bash]
# On the VPS -- run as root or in a venv
pip install andes
pip install pandapower
pip install opendssdirect

# Verify ANDES and its built-in models
python3 -c "import andes; print(andes.__version__)"
andes misc --list | grep -E "GENROU|REGCA|DFIG|Fault"
\end{lstlisting}

\subsection{Step 2: Get Standard Test Cases}

\begin{lstlisting}[style=bash]
# ANDES built-in cases (IEEE 9-bus, 14-bus, 39-bus)
python3 -c "import andes; print(andes.get_case('ieee14.json'))"
python3 -c "import andes; print(andes.get_case('ieee39.json'))"

# pandapower built-in IEEE networks
python3 -c "import pandapower.networks as pn; net = pn.case33bw(); print(net)"

# OpenDSS IEEE 13-bus feeder (comes with opendssdirect)
python3 -c "import opendssdirect as dss; print(dss.utils.get_example('13Bus'))"
\end{lstlisting}

\subsection{Step 3: Run First ANDES Transient}

\begin{lstlisting}[style=bash]
# Quick transient simulation from command line
andes run ieee9.json --routine TDS

# From Python (recommended for SAMBP integration)
python3 - <<'EOF'
import andes
ss = andes.load("ieee9.json")
ss.PFlow.run()
ss.add("Fault", {"bus": 5, "tf": 0.1, "tc": 0.25, "xf": 0.01})
ss.TDS.config.tf = 0.6
ss.TDS.run()
ss.TDS.plt.plot(ss.GENROU, "omega")   # rotor speed during fault
EOF
\end{lstlisting}

\subsection{Step 4: Wire into SAMBP Pipeline}

\begin{lstlisting}[style=bash]
# New files to create:
# 04_code/sambp/io_utils/andes_adapter.py     (see Section 4.1)
# 04_code/sambp/io_utils/opendss_adapter.py   (see Section 4.4)

# Then run sync_oc with ANDES data:
cd /root/phd_thesis
python3 04_code/sambp/sync_oc/run_milestone2.py \
        --source andes --case ieee9.json --fault-bus 5

# MC sweep with pandapower:
python3 04_code/sambp/sync_oc/main_batch_study.py \
        --source pandapower --network case33bw --n-cases 1000
\end{lstlisting}

% ============================================================
\section{What Changes vs What Stays the Same}
% ============================================================

\begin{longtable}{p{5.5cm}p{2.5cm}p{6cm}}
\toprule
\textbf{SAMBP Component} & \textbf{Status} & \textbf{Notes} \\
\midrule
\endhead
\texttt{inverse\_estimation/} LM estimator & \textbf{No change} & Consumes NumPy arrays -- source-agnostic \\[4pt]
\texttt{adaptation/} coordination logic    & \textbf{No change} & Fully decoupled from data source \\[4pt]
Stage-2 veto gate                          & \textbf{No change} & IBR flag already in FaultCase struct \\[4pt]
\texttt{models/network\_fault\_model.py}   & \textbf{Replaced} & ANDES adapter provides physics-correct waveforms \\[4pt]
\texttt{models/radial\_network.py}         & \textbf{Kept for unit tests} & Fast analytical reference for CI \\[4pt]
\texttt{config/system\_config.py}          & \textbf{Minor update} & Add ANDES case file paths \\[4pt]
\texttt{io\_utils/csv\_io.py}             & \textbf{Extend} & Add ANDES output format reader \\[4pt]
\texttt{io\_utils/andes\_adapter.py}       & \textbf{New} & $\approx$150 lines \\[4pt]
\texttt{io\_utils/opendss\_adapter.py}     & \textbf{New} & $\approx$100 lines \\[4pt]
\texttt{main\_batch\_study.py}             & \textbf{Extend} & Add pandapower source flag \\[4pt]
All TR \texttt{.tex} write-ups             & \textbf{No change} & Tool is a data source, not a model \\[4pt]
Thesis claims on LM estimator              & \textbf{Strengthened} & Higher-fidelity validation data \\
\bottomrule
\end{longtable}

% ============================================================
\section{Impact on Thesis Contributions}
% ============================================================

\subsection{Strengthened Claims}

The physics-reduction proposition and LM two-pass estimator are the
core SAMBP contributions. They are entirely independent of the simulation
tool. ANDES/OpenDSS replace the synthetic data generator, which means:

\begin{enumerate}
\item \textbf{TR-08 Monte Carlo validation} -- 4000-case sweep currently
uses synthetic Th\'evenin waveforms. Replacing with pandapower/ANDES
waveforms makes the results publishable-grade and defensible at viva.

\item \textbf{TR-56 DFIG model (87L)} -- ANDES DFIG model provides
ground-truth to validate the piecewise-ODE approximation derived
analytically in TR-56. This closes the modelling loop.

\item \textbf{TR-62 PV zone model (87L)} -- ANDES \texttt{REGCA1}
provides the current-limited waveforms against which the
2-parameter zone model is validated.

\item \textbf{TR-65 integration report} -- ANDES multi-machine transient
on IEEE 39-bus provides the full integration test that TR-65 needs to
demonstrate all nine SAMBP functions operating simultaneously.
\end{enumerate}

\subsection{Recommended Research Milestone}

\begin{center}
\begin{tabular}{lll}
\toprule
\textbf{Quarter} & \textbf{Action} & \textbf{Deliverable} \\
\midrule
Q2 2026 & Install stack, ANDES adapter for \texttt{sync\_oc}   & TR-59 updated results \\
Q3 2026 & DFIG ANDES integration for 87L (TR-56)                & TR-56 revised simulation \\
Q3 2026 & PV ANDES integration for 87L (TR-62)                  & TR-62 revised simulation \\
Q4 2026 & pandapower MC sweep replacing synthetic in TR-08       & TR-08 v2 results \\
Q1 2027 & OpenDSS for 87B distribution bus (TR-65)              & TR-65 integration test \\
Q2 2028 & IEEE 39-bus full-system test (all 9 functions)         & Chapter 7 validation \\
\bottomrule
\end{tabular}
\end{center}

% ============================================================
\section{Summary}
% ============================================================

\begin{infobox}
\textbf{Three tools -- one pip command each:}
\begin{itemize}[noitemsep]
\item \textbf{ANDES} -- transient dynamics for SG and IBR (GFL, DFIG, PV). Best fit for 87T, 87L IBR, OC, Relay 78, generator suite.
\item \textbf{OpenDSS} -- distribution-level dynamic simulation. Best fit for 87B bus differential and HIF (TR-11).
\item \textbf{pandapower} -- power flow and short-circuit Monte Carlo sweeps. Best fit for relay coordination and setting validation (TR-08, TR-46--55).
\end{itemize}
\textbf{No changes to the SAMBP core.} Only two new adapter files ($\approx$250 lines total) bridge the simulation layer to the existing LM estimator pipeline.
\end{infobox}

\vspace{1em}

The SAMBP framework is already architecturally ready for this integration.
The \texttt{io\_utils/} package, FaultCase dataclass pattern, and clean
separation between data generation and estimation mean the adapters
can be written and tested without touching any of the 136 existing
Python files.

\vfill
\begin{center}
\textcolor{gray}{\small\textit{Report generated: April 2026 $\cdot$ Anoop V Eluvathingal $\cdot$ IIT Madras}}
\end{center}

\end{document}
